\documentclass[10pt,a4paper]{amsart}
\usepackage[margin=19mm]{geometry}
\usepackage{amsmath,amssymb,mathtools}
\usepackage[scr=rsfs,cal=euler]{mathalfa}
\usepackage{xfrac}
\usepackage[unicode,hidelinks,bookmarks=false]{hyperref}
\newcommand{\ie}{i.e.\ }
\newcommand{\cf}{cf.\ }
\newcommand{\resp}{respectively\ }
\newcommand{\Subsection}{\subsection}
\providecommand{\nobreakdash}{\nobreak\hbox{-}}
\setlength{\emergencystretch}{2em}
\multlinegap0pt
\usepackage{bm}
\usepackage{bbm}
\usepackage{dsfont}
\usepackage{xspace}
\usepackage{cancel}
\usepackage{mathtools}
\usepackage{amsthm}
\makeatletter
\@ifundefined{equation}{\newtheorem{equation}{Equation}}{}
\@ifundefined{Lem}{\newtheorem{Lem}[equation]{Lemma}}{}
\makeatother
\newcommand{\N}{\mathbb{N}}
\renewcommand{\t}{\text{t}}
\title[Geometric configuration of integrally closed Noetherian doma]{Geometric configuration of integrally closed Noetherian domains}
\author{Gyu Whan Chang, Giulio Peruginelli}
\date{Source: arXiv:2601.22314v1 (CC BY 4.0)}
\begin{document}
\maketitle

\noindent\emph{Source and licence.} Geometric configuration of integrally closed Noetherian domains. Version arXiv:2601.22314v1 (CC BY 4.0), distributed under the Creative Commons Attribution 4.0 International licence, as recorded in the arXiv licence metadata for this exact version and shown on the abstract page https://arxiv.org/abs/2601.22314v1. Full rights notice: licenses/arxiv-2601.22314-CC-BY-4.0.txt.\par\smallskip

\noindent\textit{Editorial scope.} Counted unit: the Lem whose source label is \texttt{ordering} in the article (source0.tex lines 513--525, source label \texttt{ordering}), delivered together with the article's own complete original proof of it (source0.tex lines 527--554, one \texttt{proof} environment of 28 lines). No mathematics is written, repaired or re-derived: statement and proof are the article's own text line for line. Required context retained uncounted: none. Every definition, display and label of those lines is retained unchanged. Omitted from this package and not redistributed: 1--512, 526--526, 555--2416. Nothing inside a retained excerpt is omitted or altered: every retained body is delivered exactly as the article's lines are written, so no omission receipt of any kind accompanies this package. Citations: the retained excerpts carry no citation command at all, so no bibliography entry of the article is transcribed and no reference is redistributed. Reference adaptation: none was needed and none was made; every reference inside the delivered unit resolves inside it, so no unresolved reference or citation is produced. Printed slips kept literal: no printed word, symbol, hypothesis or number of the counted statement or of its proof was altered. Layout and class: the article's own class is \texttt{article} with packages including geometry, enumerate, amsmath, amssymb, amsthm, amscd, hyperref, graphicx, color, mathrsfs, xy; the wrapper is the lane's pinned \texttt{amsart} editorial wrapper at 10pt on A4, which restates the theorem-like environments the retained text needs and adds the article's own macro definitions that the retained text needs (\texttt{\textbackslash N}, \texttt{\textbackslash t}), with extra wrapper packages (bm, bbm, dsfont, xspace, cancel, mathtools). The delivered numbering is the unit's own. Assets: no retained span contains a figure environment, a graphics-inclusion command, a PDF-page inclusion, a table or a TikZ picture; no image file is copied into this package, the package declares no asset dependency and no image file is redistributed with it. Licence: version arXiv:2601.22314v1 of this article is distributed under the Creative Commons Attribution 4.0 International licence, as recorded in the arXiv licence metadata for this exact version and shown on the abstract page https://arxiv.org/abs/2601.22314v1. A full rights notice is delivered at licenses/arxiv-2601.22314-CC-BY-4.0.txt. Characters: every retained excerpt is pure ASCII, so the delivered engine, XeLaTeX with its default Latin Modern text font, reports no missing character.
\begin{Lem}\label{ordering}
Let $W_i \in \mathcal W_{\t}$ and $M_i$ be the maximal ideal of $W_i$ for $i=1,2$.
\begin{enumerate}
\item[{\em (1)}] The following statements are equivalent.
\begin{enumerate}
\item[{\em (a)}] $w_1\leq w_2$.
\item[{\em (b)}] $W_1\cap K[X]\subseteq W_2\cap K[X].$ 
\item[{\em (c)}] $M_1\cap K[X]\subseteq M_2\cap K[X]$.
\end{enumerate}
\item[{\em (2)}] $W_1=W_2$ if and only if $W_1\cap K[X]=W_2\cap K[X]$, if and only if $M_1\cap K[X] = M_2\cap K[X]$. 
\item[{\em (3)}] $(\mathcal W_{\t}, \leq)$ is a partially ordered set.
\end{enumerate}
\end{Lem} 

\begin{proof}
(1) (a) $\Rightarrow$ (c). Let $f\in M_1\cap K[X]$ be nonzero. Then $0 < w_1(f)\leq w_2(f)$, so $f\in M_2\cap K[X]$.
Thus $M_1\cap K[X]\subseteq M_2\cap K[X].$ 

(c) $\Rightarrow$ (b). Let $g\in W_1\cap K[X]$ be nonzero. If $g \in M_1$, then $g \in M_2 \cap K[X] \subseteq W_2 \cap K[X]$.
Next assume that $g \not\in M_1$ and $g \not\in W_2 \cap K[X]$. Since $W_2$ is a torsion extension of $V$,
there are $n \in \mathbb{N}$ and $b \in V$ such that $w_2(g^n) = nw_2(g) = -v(b) <0$ or $w_2(g^nb) =0$.
Note that $w_1(g^nb) = nw_1(g)+w_1(b) = v(b) >0$, so $g^nb \in M_1 \cap K[X] \subseteq M_2 \cap K[X]$
by assumption. Hence $w_2(g^nb) > 0$, a contradiction, which implies $g \in W_2 \cap K[X]$.
Thus $W_1\cap K[X]\subseteq W_2\cap K[X].$

(b) $\Rightarrow$ (a). Assume that $f\in K[X]$ is nonzero. Then either $f \in W_1$ or $\frac{1}{f} \in W_1$, and since
%{\rc [[simplified this proof by correcting a minor mistake: in general it is not true that  $V[X]\subseteq W_1\cap K[X]$]]}
%, $f(X)=\sum_i a_i X^i$. We may write $a_i=\frac{b_i}{c}$, for some $b_i,c\in V$; in particular, $f(X)=\frac{\sum_i b_i X^i}{c}$. 
%Therefore, we may restrict to polynomials $f\in V[X]$ in order to prove the inequality $w_1(f)\leq w_2(f)$. 
%Note that $V[X]\subseteq W_1\cap K[X]\subseteq W_2\cap K[X]$. In particular, for such a $f$, we have $w_1(f)\geq0$.
$W_1$ is torsion over $V$, there exist a positive integer  $n\in\N$ and $0 \neq c\in K$ such that $nw_1(f)=v(c)$. 
Hence, $w_1(\frac{f^n}{c}) = 0$. 
In particular, $\frac{f^n}{c}\in W_1\cap K[X]\subseteq W_2\cap K[X]$, so that we also have
$$0\leq w_2\left(\frac{f^n}{c}\right)\Leftrightarrow w_2(f)\geq\frac{v(c)}{n}=w_1(f),$$
so the inequality holds true. 

(2) It is clear that $W_1=W_2$ implies both of $W_1\cap K[X]=W_2\cap K[X]$ 
and $M_1\cap K[X]= M_2\cap K[X]$. Conversely, if $W_1\cap K[X]=W_2\cap K[X]$ or $M_1\cap K[X]= M_2\cap K[X]$,  
then $w_1(f)=w_2(f)$ for every $f\in K[X]$ by (1), which implies that $w_1=w_2$ on $K(X)$. Thus, $W_1=W_2$.

(3) This follows directly from (1) and (2) above.
\end{proof}

\end{document}
